\section{Expected Research Contribution}

This research will generate contributions on three levels: theoretical, methodological, and practical.

On the theoretical level, it will reveal the bidirectional coupling mechanism between neighbourhood microclimate and building energy consumption by explicitly incorporating residents' psychological perceptions and behavioural factors into the coupling framework. This addresses the current gap where human factors and feedback effects are often overlooked, and establishes a systematic framework to explain the interactions among urban thermal environment, thermal comfort, and energy consumption.

On the methodological level, this study will develop a cross-scale and integrative modelling workflow by coupling urban microclimate models (UCMs), building energy models (BEMs), and resident behavioural dynamics. The approach combines agent-based modelling for adaptive human behaviour, spatial statistical analysis to capture morphological and microclimatic heterogeneity, and path modelling to quantify causal relationships between thermal environment, perception, and energy use. Additionally, transfer learning will be employed to enable efficient adaptation of the framework to new urban contexts, while LLM-based behavioural inference will extract resident behavioural patterns from social media or survey data, reducing the need for extensive field measurements. Together, these methods establish a flexible and transferable workflow for multi-source data integration, model coupling, and scenario-based energy--comfort optimisation.

On the practical level, the research will deliver a practical computational platform to support urban planning and community-level energy optimisation strategies. It will provide optimised design guidelines for parameters such as green coverage, neighbourhood morphology, ventilation corridors, and shading configurations. Transfer learning allows rapid adaptation of the platform to new cities with minimal data, while LLM-based behavioural inference facilitates estimation of resident adaptive patterns from social media or existing surveys. Through empirical studies in Sydney and validation in other cities with similar climates, this framework will offer scientific evidence and decision-making support for advancing the green transition inhigh-density coastal cities exposed to increasing heat stress.
